\section{Experimental Setup}
\label{sec:experiments}

\subsection{Dataset}

We use an eight-family subset of the UMD Android Malware corpus,
restricted to families with at least 150 droidmon-instrumented
samples after deduplication and successful trace extraction
(Table~\ref{tab:dataset}). After a minimum-length filter of
30~API~calls, the corpus contains \textbf{9{,}337~samples}
spanning families that together cover four dominant malicious
behaviours: ad-injection (Airpush, Opfake), premium-SMS fraud
(Jisut, Genpua, SmsPay), ransomware (Fusob), and generic
trojans with dropper/rootkit behaviour (DroidKungFu, GinMaster).

\begin{table}[t]
\centering
\caption{UMD eight-family subset.}
\label{tab:dataset}
\begin{tabular}{lrr}
\toprule
Family & \# samples & Share (\%) \\
\midrule
Airpush      & 5{,}956 & 63.8 \\
DroidKungFu  & 1{,}315 & 14.1 \\
Opfake       &   615 & 6.6  \\
Jisut        &   534 & 5.7  \\
GinMaster    &   529 & 5.7  \\
Genpua       &   313 & 3.4  \\
SmsPay       &   174 & 1.9  \\
Fusob        &   166 & 1.8  \\
\midrule
\textbf{Total} & \textbf{9{,}337} & 100.0 \\
\bottomrule
\end{tabular}
\end{table}

The imbalance ratio between the largest and smallest family is
approximately $34\times$, which we address through class-weighted
cross-entropy training (Section~\ref{sec:training}) and
macro-averaged reporting.

\subsection{Baselines}

We compare GAME-Mal against six baselines, all trained and
evaluated under the identical fold splits described below:

\begin{itemize}
  \item \textbf{Random Forest} with $n=200$ trees over the
    pruned-rule feature vectors.
  \item \textbf{Linear SVM} with one-vs-rest, $C=1.0$, balanced
    class weights.
  \item \textbf{Decision Tree} with default scikit-learn settings.
  \item \textbf{Gaussian Naïve Bayes} over the rule features.
  \item \textbf{MarkovPruning}: our faithful reproduction of
    D'Angelo et al.~\cite{dangelo2023association}, classifying via
    Eq.~\ref{eq:rho}. This serves as our primary statistical
    baseline.
  \item \textbf{Plain Transformer}: an ablation of GAME-Mal with
    the sigmoid gates removed (equivalent to standard multi-head
    attention) but otherwise identical architecture, training
    recipe, and hyperparameters.
\end{itemize}

\subsection{Evaluation protocol}

We use \textbf{stratified 3-fold cross-validation}. Each family's
samples are partitioned into three disjoint folds, preserving the
family's global proportion. For each fold, the other two folds
are used for training and rule mining, the held-out fold for
evaluation. Reported numbers are means over the three folds.

We report \textbf{accuracy}, \textbf{macro-averaged F1-score},
and \textbf{macro-averaged one-vs-rest AUROC}. Macro-averaging
prevents the majority class (Airpush at 63.8\%) from dominating
the headline numbers. All deep models are evaluated with the
checkpoint that achieved the highest validation macro-F1 during
training, not the last-epoch checkpoint.

\subsection{Implementation}

All code is implemented in Python 3.9 with PyTorch 2.2 and
scikit-learn 1.4. Experiments run on an Apple~M-series Mac,
using the PyTorch MPS backend for GPU acceleration. A full
three-fold run (seven classifiers per fold, including training
GAME-Mal for up to 40 epochs) completes in approximately
2.5~hours. Our full training and evaluation pipeline, together
with pre-computed vocabularies, Markov rules, and trained model
weights, is released at \url{https://github.com/<redacted>/game-mal}.
